\documentclass{article}
\usepackage[T1]{fontenc}
\usepackage[utf8]{inputenc}

% Linespread command allows you to change line spacing for the entire document
\linespread{1.18}

% Tweak page margins
\addtolength{\oddsidemargin}{-.875in}
\addtolength{\evensidemargin}{-.875in}
\addtolength{\textwidth}{1.75in}

\addtolength{\topmargin}{-.875in}
\addtolength{\textheight}{1.75in}

\usepackage{natbib}
\usepackage{hyperref}
\usepackage{xcolor}
\usepackage{xspace}
\usepackage{fancyhdr}
\hypersetup{
    colorlinks,
    linkcolor={red!50!black},
    citecolor={blue!50!black},
    urlcolor={blue!80!black}
}

\newcommand{\HRule}{\rule{\linewidth}{0.5mm}}
\newcommand{\Hrule}{\rule{\linewidth}{0.3mm}}

% Project specific macros
\newcommand{\graphite}{GRAPHITE\xspace}
\newcommand{\wave}{WAVE\xspace}

% School specific macros
\newcommand{\schoolShort}{Stanford\xspace}
\newcommand{\school}{Stanford\xspace}
\newcommand{\schoolLong}{\xspace}

\newcommand{\profOne}{Prof. Matei Zaharia\xspace}
\newcommand{\profTwo}{Prof. Christos Kozyrakis\xspace}
\newcommand{\profThree}{Prof. Jure Leskovec\xspace}

% Creates header for each page
\usepackage{fancyhdr}
\pagestyle{fancy}
\fancyhf{}
\fancyhead[LE,RO]{\header\hskip\linepagesep\vfootline\thepage}
\newskip\linepagesep \linepagesep 5pt\relax
\def\vfootline{%
    \begingroup
    	\rule[-10pt]{0.75pt}{25pt}
    \endgroup
}
\def\header{%
	\begin{minipage}[]{120pt}
		\hfill Arunima Chowdhury 		% Applicant Name
    	\par \hfill 				% Formatting boilerplate
    	CS, PhD, Fall 2027 			% Area, Program, Cycle, Year
    \end{minipage}
}
\fancyhead[RE,LO]{Statement of Purpose | \schoolLong}
\renewcommand\headrulewidth{0pt}

\begin{document}

%% Why do you wish to attend graduate school? What would you like to study? Keep it broad, details come-in later
\textbf{Introduction:} My research interests center on developing reliable supervision for foundation models. A recent shift has shown how powerful it can be: instead of learning only from human-written examples, models can now be trained with rewards that check whether their answers are correct, and through this process they learn to reason step by step. As training continues, they even learn on their own to think longer before responding, which opens a second way to improve them: spending more computation at test time. Yet this approach works best where answers can be checked automatically, such as math and code. In many other settings, models must rely on learned reward models, which judge outputs less reliably and can be exploited during training. In my PhD, I hope to contribute to this area by studying how learned reward models can give more reliable, step-level feedback, both for training models and for selecting among their outputs at test time.



%% Describe 2-3 past projects that might be relevant to your research interests. (10-12 lines per project)

% PROJECT 1: med prm

\textbf{Research Experience:} I first encountered this challenge during my master’s research on multimodal medical reasoning. While studying process reward models, I became interested in a key limitation: they can identify incorrect reasoning steps, but often provide little insight into what caused the error. For medical image analysis, a model must locate the correct region, recognize what it contains, and reason from it to a diagnosis. A model might misinterpret an image or draw an incorrect clinical conclusion from accurately identified visual evidence. Although both lead to incorrect reasoning, they represent different types of errors that require different forms of supervision.

To address this, I am working with Prof. Ruihong Huang at Texas A\&M on a typed process reward model that judges each reasoning step on two separate axes: whether its visual claims are correct, and whether its clinical inference holds given those claims. To build supervision without manually labeling each reasoning step, I sample reasoning chains from Qwen2.5-VL-7B, parse them into visual claims and inferences, check visual claims against available expert annotations, and evaluate clinical inferences against retrieved medical knowledge. If an inference relies on a visual claim that turns out to be wrong, I label it as dependent rather than counting it as a separate error. To ensure the model is actually using the image, I test whether its visual judgments retain when the image is replaced with a blank one. I plan to use the trained model both for Best-of-N selection among candidate reasoning chains and as a step-level reward in GRPO to train models that reason from the image rather than from language priors.

% PROJECT 2: bangla

I encountered a related challenge in an earlier project on Bangla speech emotion recognition, where labeled emotional speech is scarce. After fine-tuning Whisper and HuBERT on a small labeled dataset, I found that performance varied sharply across emotion classes, especially those with few examples. To make use of a larger unlabeled corpus, I implemented iterative pseudo-labeling, where I added the models’ high-confidence predictions to subsequent training rounds. But confidence alone was unreliable: confidently wrong predictions reinforced existing class biases, particularly between acoustically similar emotions. To tackle this problem, I combined confidence filtering with class-balanced selection. This raised macro-F1 from 61.8\% to 68.4\%, with the largest gains in underrepresented classes. The experience taught me that adding more data is not enough if the supervision itself is unreliable. What matters is whether we can identify and control errors in the labels we learn from, a concern I now revisit when constructing step-level labels for reasoning.

% PROJECT 3: Dance Motion Feedback System

My undergraduate research at Purdue gave me an early appreciation for how much the quality of model-generated feedback depends on the evidence available to it. I spent two years building an automated feedback system for modern dance that compares student performances against instructor references and generates personalized corrective guidance. The project started with 3D motion tracking for live theatre, where I needed fast, reliable pose estimation for multiple dancers. YOLO-Pose with stereo depth sensing gave me that --- roughly 3x faster and 70\% more accurate than TensorFlow PoseNet. I classified movements across 19 categories using a fine-tuned SkateFormer and aligned student and instructor performances using FastDTW to measure joint-level deviations. To turn these measurements into actionable coaching, I used Gemini as a feedback generator. I provided it with structured, joint-level movement deviations and prompted it to generate corrective feedback tied to those measurements, rather than relying on the model to identify movement errors on its own. The work received the Best Research Talk Award at Purdue’s Undergraduate Research Conference.

%% Non-research accomplishments (e.g. Grades, Academic Service, Work experience) (10-12 lines)
Alongside my research, my graduate coursework at Texas A\&M University has strengthened my foundation in machine learning and reinforcement learning. I have maintained a 3.9 GPA while taking courses in Deep Reinforcement Learning and Machine Learning and LLMs. In Deep Reinforcement Learning, a teammate and I investigated how reinforcement learning could optimize workforce allocation in a call center, where employees must balance serving incoming customers with completing back-office tasks. We formulated the problem as a sequential decision-making process, with an agent allocating each employee's time at 15-minute intervals throughout a simulated workday. To model the environment, we developed a discrete-event simulation incorporating stochastic customer arrivals, correlated service and patience times, and customer abandonment. We then trained a Deep Q-Network (DQN) using a reward function that penalized customer waiting times, abandonment, employee overload, and unfinished back-office work. Across 200 simulated days, the learned policy substantially reduced waiting times and abandonment compared with random assignment and distributed back-office responsibilities more evenly. However, it completed fewer back-office tasks, whose penalty was assessed only at the end of the day. This tradeoff raised an important question about how reward design and delayed feedback influence an agent's learned behavior, particularly when optimizing competing objectives over long decision sequences. The experience strengthened my interest in reinforcement learning, especially credit assignment and the challenges of learning effective policies from delayed rewards.

Beyond research and coursework, I have developed teaching, mentoring, and leadership experience. As a teaching assistant for Purdue's Data Mine program in collaboration with Merck, I led a team of nine student developers, coordinating project milestones, assigning development tasks, reviewing code, and helping teammates troubleshoot technical challenges. I also served as a liaison between the student team and Merck, communicating project progress and ensuring that our implementation aligned with the project's objectives. At Texas A\&M, I have twice served as a grader for the Computer Engineering senior capstone and currently grade for Software Engineering, evaluating student projects against established requirements and assessment criteria. At Purdue, I was recognized on the Dean's List and served as Vice President of the Bangladesh Student Association, where I secured \$47,000 in funding from the Student Fee Advisory Board to support cultural events for the Bangladeshi community. My professional experience also includes a data analyst internship at Robi Axiata, where I developed forecasting models using more than 100,000 subscriber records and built dashboards to support business decision-making. These experiences have strengthened my ability to lead collaborative projects, mentor other students, and communicate technical ideas across academic and professional settings.

% Grades

% TA and Academic Service

% Industry

%% Why this school? List professors you would like to work with and why? (10-12 Lines)



%% Summary (3-4 Lines)


% Add some blank space between text and references
\vspace{0.125in}

% References

% **NOTE**: There are better ways to manage citations in LaTeX, most notably using a bibTeX. I wanted to have greater control on how citations were spaced and formatted and therefore ended up hardcoding them here. Your mileage may wary!



\end{document}

% That's All Folks.

% Best of luck, you got this! :)
